\chapter{Generating Functions and Transforms}\label{ch:generating-functions}

\chapterinfo{STAT2003}{\cref{ch:random-variables,ch:expectation-moments}}

% Sources: STAT2003 handwritten notes (pp. 26-27); STAT2003 course notes, Section 2.4.

Many calculations and manipulations involving probability distributions are made easier by transforms. This chapter covers two of them: the probability generating function, for non-negative integer-valued random variables, and the moment generating function. Both determine the distribution uniquely, and both give the moments of a random variable by differentiation. They are particularly useful for sums of independent random variables (\cref{sec:sums-transforms}).

% ==========================================================================
\section{Probability generating functions}\label{sec:pgf}
% ==========================================================================

\begin{definition}[Probability generating function]\label{def:pgf}
Let $X$ be a non-negative, integer-valued random variable. The \emph{probability generating function} (PGF) of $X$ is the function $G \colon [0,1] \to [0,1]$ defined by
\[
  G(z) \coloneqq \E(z^X) = \sum_{x=0}^\infty z^x \, \Prob(X = x).
\]
\end{definition}

\begin{example}[PGF of the Poisson distribution]\label{ex:pgf-poisson}
Let $X$ have pmf
\[
  f(x) = e^{-\lambda} \frac{\lambda^x}{x!}, \qquad x = 0, 1, 2, \ldots
\]
(the Poisson distribution, \cref{ch:common-distributions}). Its PGF is
\[
  G(z) = \sum_{x=0}^\infty z^x e^{-\lambda} \frac{\lambda^x}{x!} = e^{-\lambda} \sum_{x=0}^\infty \frac{(\lambda z)^x}{x!} = e^{-\lambda} e^{\lambda z} = e^{-\lambda(1 - z)}.
\]
\end{example}

\begin{theorem}[The pmf from the PGF]\label{thm:pgf-pmf}
If $X$ has PGF $G$, then
\[
  \Prob(X = x) = \frac{1}{x!} \frac{\dd^x}{\dd z^x} G(z) \Big|_{z=0}, \qquad x = 0, 1, 2, \ldots
\]
\end{theorem}

\begin{proof}
By definition,
\[
  G(z) = z^0 \Prob(X = 0) + z^1 \Prob(X = 1) + z^2 \Prob(X = 2) + \cdots.
\]
Substituting $z = 0$ gives $G(0) = \Prob(X = 0)$. Differentiating once,
\[
  G'(z) = \Prob(X = 1) + 2z \Prob(X = 2) + 3z^2 \Prob(X = 3) + \cdots,
\]
so $G'(0) = \Prob(X = 1)$. Differentiating again gives $G''(0) = 2 \Prob(X = 2)$, and in general the $x$th derivative of $G$ at zero is $G^{(x)}(0) = x! \, \Prob(X = x)$.
\end{proof}

So the PGF has the \emph{uniqueness property}: two distributions are the same if and only if their PGFs are the same. The moments of $X$ follow from the derivatives of $G$ at $z = 1$.

\begin{theorem}[Moments from the PGF]\label{thm:pgf-moments}
If $X$ has PGF $G$, then
\[
  \E(X) = G'(1), \qquad \E(X(X-1)) = G''(1), \qquad \E(X^k) = \Bigl(z \frac{\dd}{\dd z}\Bigr)^k G(z) \Big|_{z = 1^-}.
\]
Consequently,
\[
  \Var(X) = G''(1) + G'(1) - \bigl(G'(1)\bigr)^2.
\]
\end{theorem}

\begin{proof}
Differentiating $G(z) = \E(z^X)$ with respect to $z$ gives
\begin{align*}
  G'(z) &= \E\bigl(X z^{X-1}\bigr), \\
  G''(z) &= \E\bigl(X(X-1) z^{X-2}\bigr), \\
  G'''(z) &= \E\bigl(X(X-1)(X-2) z^{X-3}\bigr),
\end{align*}
and so on (write the expectation as a sum and differentiate term by term; some care is needed with infinite sums). Setting $z = 1$ gives $\E(X)$ and $\E(X(X-1))$. Since $\E(X^2) = \E(X(X-1)) + \E(X)$,
\[
  \Var(X) = \E(X^2) - (\E X)^2 = G''(1) + G'(1) - \bigl(G'(1)\bigr)^2.
\]
The general formula for $\E(X^k)$ is not proved in this course.
\end{proof}

% ==========================================================================
\section{Moment generating functions}\label{sec:mgf}
% ==========================================================================

\begin{definition}[Moment generating function]\label{def:mgf}
The \emph{moment generating function} (MGF) of a random variable $X$ is the function $M \colon I \to [0, \infty)$ given by
\[
  M(s) = \E(e^{sX}) = \begin{cases}
    \displaystyle \sum_x e^{sx} f(x) & \text{if $X$ is discrete with pmf $f$}, \\[10pt]
    \displaystyle \int_{-\infty}^\infty e^{sx} f(x) \dd x & \text{if $X$ is continuous with pdf $f$}.
  \end{cases}
\]
Here $I$ is an open interval containing $0$ on which the sum or integral is finite for all $s \in I$.
\end{definition}

The MGF also has the uniqueness property: two distributions are the same if and only if their MGFs are the same. As the name suggests, the derivatives of $M$ at $s = 0$ give the moments of $X$.

\begin{theorem}[Moments from the MGF]\label{thm:mgf-moments}
If $\E(X^n)$ exists, then $M$ has $n$ derivatives, and
\[
  \E(X^n) = M^{(n)}(0).
\]
Consequently,
\[
  \Var(X) = M''(0) - \bigl(M'(0)\bigr)^2.
\]
\end{theorem}
